\chapter{Mathematical Writing}\label{app:mathematical-writing}

The guidelines below summarize how to write a clear mathematical thesis.
They are taken from classic sources on mathematical and scientific writing,
which are worth reading in full:
Halmos~\cite{halmos1970HowWriteMathematics},
Knuth et al.~\cite{knuth1989MathematicalWriting},
Gopen and Swan~\cite{gopen1990ScienceScientificWriting},
Lamport~\cite{lamport2012HowWrite21st},
Zobel~\cite{zobel2014WritingComputerScience}, and
Higham~\cite{higham2020HandbookWritingMathematical}.
The guiding principle behind all of them is that you write for the
\emph{reader}: a fellow student who knows the basics of computer science
but not your topic.
Everything that makes the reader stop, search, or guess is a problem of the
text, not of the reader.

%-----------------------------------------------------------------------
\section{Structure and Flow}\label{sec:mw-structure}
%-----------------------------------------------------------------------

\begin{itemize}
    \item \textbf{Intuition first, then formality, then an example.}
          Before each definition or theorem, say in words what it formalizes
          and why it is needed.
          After it, illustrate it with an example.
          Use one running example throughout a chapter instead of many
          unrelated ones.
    \item \textbf{One topic per paragraph.}
          The first sentence of a paragraph states its topic, the remaining
          sentences develop it.
    \item \textbf{Old information first, new information last.}
          Start a sentence with what the reader already knows and put the new
          or important part at the end, where the reader expects the emphasis.
          Keep the subject and its verb close together.
    \item \textbf{Signposting.}
          Begin every chapter by saying what it contains and end it by
          summarizing what was achieved and where it is used next.
          Within long sections, tell the reader where the argument is going
          (``To prove \cref{thm:main}, we first show two auxiliary lemmas.'').
    \item \textbf{Present only what you need.}
          The preliminaries contain exactly the notions used later.
          Do not include definitions or results ``for completeness'' that
          play no role in your thesis.
\end{itemize}

%-----------------------------------------------------------------------
\section{Definitions and Notation}\label{sec:mw-notation}
%-----------------------------------------------------------------------

\begin{itemize}
    \item \textbf{Define everything before you use it,} and define it only
          once.
          If a notion was defined many pages before, recall it briefly and
          refer to the definition (``Recall that a term is \emph{linear} if
          \ldots{} (see \cref{def:first-notion}).'').
    \item \textbf{Emphasize the defined term} where it is defined, e.g.,
          ``A relation is called \emph{terminating} if \ldots''.
    \item \textbf{Definitions use ``if'' and ``$=$''.}
          In a definition, ``if'' already means ``if and only if'', so never
          write ``iff'' or ``if and only if'' in a definition.
          Likewise, define objects with ``$=$'' (``Let $f(x) = x + 1$.''),
          not with \lstinline|:=|, \lstinline|\coloneqq|, or
          \lstinline|::=|.
          Use ``iff'' or ``if and only if'' only in statements that you prove,
          e.g., in a theorem or lemma, where both directions have to be shown.
    \item \textbf{One symbol, one meaning.}
          Never use the same symbol for two different objects, and never use
          two symbols for the same object.
          Keep the conventions consistent: e.g., lowercase letters for
          elements, uppercase letters for sets, calligraphic letters for
          systems.
    \item \textbf{Do not introduce notation that is used only once.}
          Every new symbol is a burden for the reader; a few words are often
          clearer.
    \item \textbf{One word, one concept.}
          In technical writing, varying the words for the same concept (term,
          expression, object) suggests that they mean different things.
          Pick one word and stick to it.
    \item \textbf{Use standard notation} wherever it exists, e.g., the
          notation of the standard textbook of your area.
          If your thesis uses many symbols, add a list of notation to the
          front matter.
\end{itemize}

%-----------------------------------------------------------------------
\section{Formulas in the Text}\label{sec:mw-formulas}
%-----------------------------------------------------------------------

Mathematical formulas are part of the sentences they appear in and follow the
same rules of grammar and punctuation.

\begin{itemize}
    \item \textbf{Do not start a sentence with a symbol.}
    \item \textbf{Do not separate two formulas by only a punctuation mark};
          put words in between.
    \item \textbf{Do not use logical symbols as abbreviations in prose.}
          Symbols like $\forall$, $\exists$, $\Rightarrow$, or $\wedge$
          belong inside formulas, not between words.
    \item \textbf{Punctuate displayed formulas.}
          A displayed formula ends with a comma or a period if the sentence
          does, e.g.,
          \begin{equation}\label{eq:mw-geometric-sum}
              \sum_{i=0}^{n} 2^i = 2^{n+1} - 1,
          \end{equation}
          where $n \in \IN$.
    \item \textbf{Numbering.}
          Use \lstinline|\begin{equation}| with a \lstinline|\label| for
          displayed formulas that you refer to, and refer to them with
          \lstinline|\cref|, e.g., \lstinline|\cref{eq:mw-geometric-sum}|
          yields \cref{eq:mw-geometric-sum}.
          Use \lstinline|\begin{equation*}| for formulas that are never
          referenced.
    \item \textbf{Words in formulas} are set with \lstinline|\text|, and
          operator names with \lstinline|\operatorname| or the predefined
          commands such as \lstinline|\max| and \lstinline|\log|, so that they
          are not typeset as products of variables:
          $\max(x, y)$ instead of $max(x, y)$.
    \item \textbf{Use the right symbols:}
          \lstinline|\colon| for the type of a function
          ($f \colon \IN \to \IN$),
          \lstinline|\mid| in set comprehensions
          ($\set{x \in \IN \mid x > 0}$),
          \lstinline|\ldots| between commas ($x_1, \ldots, x_n$), and
          \lstinline|\cdots| between operators ($x_1 + \cdots + x_n$).
\end{itemize}

The following table shows typical mistakes and how to avoid them.

\begin{table}[H]
    \centering
    \caption{Typical mistakes in mathematical prose.}
    \label{tab:mw-mistakes}
    \begin{tabular}{p{0.44\textwidth}p{0.48\textwidth}}
        \toprule
        \textbf{Avoid} & \textbf{Prefer} \\
        \midrule
        $x$ is positive, so \ldots
        & The value $x$ is positive, so \ldots \\
        \addlinespace
        Since $x > 0$, $y < 0$.
        & Since $x > 0$, we have $y < 0$. \\
        \addlinespace
        $\forall x \in S$ we have $x > 0$.
        & For all $x \in S$, we have $x > 0$. \\
        \addlinespace
        If $n$ is even $\Rightarrow$ $n^2$ is even.
        & If $n$ is even, then $n^2$ is even. \\
        \addlinespace
        As shown in Theorem 2.2.2, \ldots
        & As shown in \cref{thm:key-theorem}, \ldots \\
        \addlinespace
        This follows from it.
        & This lemma follows from \cref{thm:key-theorem}. \\
        \addlinespace
        Obviously, the relation is terminating.
        & The relation is terminating, since every rule decreases the size
          of the term. \\
        \addlinespace
        (Definition) A term is \emph{linear} iff \ldots
        & (Definition) A term is \emph{linear} if \ldots \\
        \addlinespace
        (Definition) Let $f(x) \coloneqq x + 1$.
        & (Definition) Let $f(x) = x + 1$. \\
        \bottomrule
    \end{tabular}
\end{table}

%-----------------------------------------------------------------------
\section{Theorems and Proofs}\label{sec:mw-proofs}
%-----------------------------------------------------------------------

\begin{itemize}
    \item \textbf{Self-contained statements.}
          A theorem lists all of its assumptions.
          Avoid statements like ``Under the assumptions above, \ldots'', since
          the reader has to search for what ``above'' means.
    \item \textbf{Explain the significance} of every result: why it is
          needed, how it is used later, and what it means intuitively.
    \item \textbf{Announce the structure of a proof} at its beginning
          (``We prove the claim by induction on $n$.'',
          ``Assume for a contradiction that \ldots'').
          Mark every case clearly and state the induction hypothesis
          explicitly.
    \item \textbf{Justify every step} and refer to the definition or result
          it uses by its label.
          Structure long proofs hierarchically into numbered steps, each with
          its own short justification.
    \item \textbf{Avoid ``obviously'', ``clearly'', and ``trivially''.}
          If a step is obvious, these words are unnecessary; if it is not, they
          hide a gap.
          Give the reason instead.
    \item \textbf{Long proofs go to the appendix.}
          Give a proof sketch in the main text and move the full proof to
          \cref{app:omitted-proofs}.
    \item \textbf{Apply the result to an example} after the proof, ideally
          to the running example.
\end{itemize}

The following lemma and its proof illustrate these guidelines.
The lemma explains its significance before the statement, the proof
announces its structure, marks both cases, states the induction hypothesis,
and justifies the only non-trivial step.

We now show that the sum of the first $n$ odd numbers is a square number.
This allows us to compute such sums without adding up all summands.

\begin{Lemma}{Sum of Odd Numbers}{mw-odd-sum}
    For all $n \in \IN$, we have
    \begin{equation}\label{eq:mw-odd-sum}
        \sum_{i=1}^{n} (2i - 1) = n^2.
    \end{equation}
\end{Lemma}

\begin{proof}
    We prove the claim by induction on $n$.

    \textbf{Base case ($n = 0$).}
    The sum on the left-hand side of \cref{eq:mw-odd-sum} is empty, so it is
    $0 = 0^2$.

    \textbf{Induction step ($n \to n + 1$).}
    As induction hypothesis, assume that \cref{eq:mw-odd-sum} holds for $n$.
    By splitting off the last summand and applying the induction hypothesis,
    we obtain
    \begin{equation*}
        \sum_{i=1}^{n+1} (2i - 1)
        = \sum_{i=1}^{n} (2i - 1) + 2(n + 1) - 1
        = n^2 + 2n + 1
        = (n + 1)^2,
    \end{equation*}
    so \cref{eq:mw-odd-sum} also holds for $n + 1$.
\end{proof}

%-----------------------------------------------------------------------
\section{Language and Style}\label{sec:mw-language}
%-----------------------------------------------------------------------

\begin{itemize}
    \item \textbf{Use ``we''} consistently, meaning the author and the reader
          (``We now show that \ldots'').
          Never use ``I''!
    \item \textbf{Prefer the active voice and the present tense}
          (``\Cref{thm:main} shows \ldots'').
          Use the past tense only for things that happened at a specific
          time, e.g., when you describe the experiments you ran.
    \item \textbf{Be precise.}
          Replace vague words like ``very'', ``quite'', ``some'', or ``a lot''
          by concrete statements or numbers
          (``solves 23 more examples'' instead of ``solves many more
          examples'').
          Avoid ambiguous ``it'' and ``this'' and name the object instead.
    \item \textbf{Abbreviations} are spelled out at their first use, e.g.,
          ``dependency pairs (DPs)'', and are not used in the abstract.
    \item \textbf{Spelling.}
          Use either American or British English consistently, and run a
          spell checker before submission.
          Write ``e.g.,'' and ``i.e.,'' with a comma afterwards.
\end{itemize}

%-----------------------------------------------------------------------
\section{Figures and Tables}\label{sec:mw-figures}
%-----------------------------------------------------------------------

\begin{itemize}
    \item \textbf{Every figure and table is referenced} in the text with
          \lstinline|\cref|, and the text explains what the reader should
          see in it.
    \item \textbf{Captions are self-contained}: a reader who only looks at
          the figure and its caption should understand what is shown,
          including the meaning of all abbreviations and units.
    \item \textbf{Tables} use horizontal rules only (\lstinline|\toprule|,
          \lstinline|\midrule|, \lstinline|\bottomrule| from
          \lstinline|booktabs|) and no vertical rules, see
          \cref{tab:mw-mistakes} and \cref{tab:results}.
    \item \textbf{Colors} must not be the only way to distinguish
          information: the thesis has to be readable in black and white and
          by readers with color vision deficiency.
\end{itemize}
